% !TeX program = xelatex
\documentclass[12pt, a4paper]{article}

% Packages:
\usepackage[
    ignoreheadfoot,
    top=1.25cm,
    bottom=2cm,
    left=1.25cm,
    right=1.25cm,
    footskip=1.0 cm
]{geometry}
\usepackage[explicit]{titlesec}
\usepackage{xcolor}
\definecolor{primaryColor}{RGB}{0, 79, 144}
\usepackage{enumitem}
\usepackage{fontawesome5}
\usepackage[
    pdftitle={Сергей Ермолаев CV},
    pdfauthor={Сергей Ермолаев},
    pdfcreator={LaTeX with RenderCV},
    colorlinks=true,
    urlcolor=primaryColor
]{hyperref}
\usepackage{calc}
\usepackage{bookmark}
\usepackage{changepage}
\usepackage{paracol}
\usepackage{ifthen}
\usepackage{needspace}
\usepackage{etoolbox} % for \patchcmd


\usepackage{fontspec}
\usepackage{polyglossia}
\setmainlanguage{russian}
\setotherlanguage{english}
\defaultfontfeatures{Ligatures=TeX, Scale=MatchLowercase}

\setmainfont{Inter_18pt}[
  Extension = .ttf,
  UprightFont = *-Regular,
  BoldFont = *-Bold,
  ItalicFont = *-Italic,
  BoldItalicFont = *-BoldItalic
]


\newfontfamily\cyrillicfont{Inter_18pt}[
  Extension = .ttf,
  UprightFont = *-Regular,
  BoldFont = *-Bold,
  ItalicFont = *-Italic
]

\newfontfamily\cyrillicfontsf{Inter_18pt}[
  Extension = .ttf,
  UprightFont = *-Regular,
  BoldFont = *-Bold,
  Script=Cyrillic
]


% Some settings:
\AtBeginEnvironment{adjustwidth}{\partopsep0pt}
\pagestyle{plain}
\setcounter{secnumdepth}{0}
\setlength{\parindent}{0pt}
\setlength{\topskip}{0pt}
\setlength{\columnsep}{0.15cm}

\titleformat{\section}{
  \needspace{4\baselineskip}
  \Large\color{primaryColor}
}{}{}{\textbf{#1}\hspace{0.1cm}\titlerule[0.8pt]\hspace{-0.1cm}}[]

\titleformat{\subsection}{
}{}{}{\textit{#1}}[]

\titlespacing{\section}{-1pt}{0.08 cm}{0.05 cm}
\titlespacing{\subsection}{-1pt}{0.08 cm}{0.05 cm}
\newenvironment{highlights}{
  \begin{itemize}[
    topsep=0.05 cm,
    parsep=0.10 cm,
    partopsep=0pt,
    itemsep=0pt,
    leftmargin=0.4 cm + 10pt
  ]
}{
  \end{itemize}
}

\newenvironment{onecolentry}{
  \begin{adjustwidth}{0.0cm}{0.1cm}
}{
  \end{adjustwidth}
}

\newenvironment{twocolentry}[2][]{
  \onecolentry
  \def\secondColumn{#2}
  \setcolumnwidth{\fill, 3.5 cm}
  \begin{paracol}{2}
}{
  \switchcolumn \raggedleft \secondColumn
  \end{paracol}
  \endonecolentry
}


\newenvironment{header}{
  \setlength{\topsep}{0pt}\par\kern\topsep\centering\color{primaryColor}\linespread{1.5}
}{
  \par\kern\topsep
}

\newcommand{\eduentry}[4]{% #1 вуз; #2 локация; #3 степень/направление; #4 даты
  \noindent
  \begin{tabular*}{\linewidth}{@{}l@{\extracolsep{\fill}}r@{}}
    \textbf{#1} & #2\\[-0.2em]
    \textit{#3} & \textit{#4}\\
  \end{tabular*}\par
}

\newcommand{\subeduentry}[2]{% #1 название; #2 даты
  \noindent
  \begin{tabular*}{\linewidth}{@{}l@{\extracolsep{\fill}}r@{}}
    {#1} & \textit{#2}
  \end{tabular*}\par
}


% save the original href command in a new command:
\let\hrefWithoutArrow\href
% new command for external links:
\renewcommand{\href}[2]{\hrefWithoutArrow{#1}{\ifthenelse{\equal{#2}{}}{ }{#2 }\raisebox{.15ex}{\footnotesize \faExternalLink*}}}

\begin{document}

\newcommand{\AND}{\unskip
  \cleaders\copy\ANDbox\hskip\wd\ANDbox
  \ignorespaces
}
\newsavebox\ANDbox
\sbox\ANDbox{}

% \placelastupdatedtext
\begin{header}
  {\fontsize{28pt}{34pt}\selectfont \textbf{Сергей Ермолаев}}

  \normalsize
  \mbox{\hrefWithoutArrow{https://t.me/j1nsei/}{{\footnotesize\faTelegramPlane}\hspace*{0.13cm}}}
  \kern 0.25 cm
  \mbox{\hrefWithoutArrow{https://www.linkedin.com/in/saermolaev/?locale=ru-RU}{{\footnotesize\faLinkedin}\hspace*{0.13cm}}}
  \kern 0.25 cm
  \AND%
  \kern 0.25 cm%
  \mbox{\hrefWithoutArrow{mailto:saermolaev.w@gmail.com}{{\footnotesize\faEnvelope[regular]}\hspace*{0.13cm}saermolaev.w@gmail.com}}%
  \kern 0.25 cm%
  \AND%
  \kern 0.25 cm%
  \mbox{\hrefWithoutArrow{tel:+7 (930) 802-31-25}{{\footnotesize\faPhone*}\hspace*{0.13cm}+7 (930) 802-31-25}}%
  \kern 0.25 cm%
  \AND%
  \kern 0.25 cm%
  \mbox{{\footnotesize\faMapMarker*}\hspace*{0.13cm}Нижний Новгород}%
  \kern 0.25 cm%
  \AND%
  \kern 0.25 cm%
  \mbox{\hrefWithoutArrow{https://j1nsei.github.io/cv/en.html?utm_source=ru_pdf}{{\footnotesize\faLanguage}\hspace*{0.13cm}ENG CV.}}
\end{header}

\section{Обо мне}

\begin{onecolentry}
    {\sloppy
    \textbf{AI / Machine Learning Engineer с 3 годами коммерческого опыта.} Работал с классификацией, компьютерным зрением и обработкой речи. Сейчас основной фокус — \textbf{RAG, LLM-инференс, AI-агенты и гибридный поиск}. Сильная сторона — быстрое прототипирование ML/AI-решений, от проверки гипотез до готового демо. Довожу решения до production с оценкой качества и мониторингом, оптимизирую их под требования к скорости и вычислительным ресурсам.
    \par\fussy}
\end{onecolentry}


\section{Опыт}

\begin{twocolentry}{Нижний Новгород \\ Нояб. 2023 -- Н. в.}
  \textbf{Протон} \textit{(IT‑компания ГК \textbf{"АГАТ"}, \href{https://www.forbes.ru/biznes/535642-vyrucka-25-krupnejsih-avtodilerov-rossii-vyrosla-do-2-4-trln-rublej}{топ-5 автодилеров РФ}, 120+ автосалонов)}\newline
  ML Engineer:
  \begin{highlights}
    \item \textbf{Многометочная классификация звонков}: F1-samples \textbf{79\% → 94\%}, EMR \textbf{3\% → 61.6\%}; инференс на CPU за доли секунды, \textbf{>1000 звонков ежедневно}. Подготовка данных, транскрибация, стратификация, эксперименты и интеграция в production. По оценке аналитиков компании, экономический эффект — \textbf{\textasciitilde170 млн ₽}.
    \item \textbf{RAG‑ассистент (RAGFlow + vLLM)}: user satisfaction \textbf{4.7/5}, faithfulness \textbf{95\%}, \textbf{TTFT p99 4.6s} под нагрузкой с большим контекстом. Подбор моделей, доработка RAGFlow под корпоративную БЗ, собственный пайплайн исправления и парсинга документов. Развёртывание в production, мониторинг, трейсинг и оценка качества.
    \item \textbf{Агент} маршрутизации тикетов в системе управления бизнесом GANDIVA с использованием RAG, Guardrails, Memory. Полный цикл разработки, включая мониторинг и оценку качества / производительности \textit{(LangGraph, RAGFlow, LangFuse, vLLM)}.
    \item \textbf{Система подбора автозапчастей}: снижение обращений к подборщикам с \textbf{300 до 20 в день (−93\%)}. Реализация гибридного поиска с собственной функцией ранжирования и порогами отбора, создание словаря сокращений. Развёртывание сервиса на \textbf{FastAPI + Docker} в production.
    \item \textbf{LLM‑разметка} звонков (Mistral/Qwen, zero-/few-shot, \textbf{LoRA Qwen2.5‑3B}): Accuracy до 90\% (LoRA), прод — XLinear 88\% (CPU онлайн/офлайн).
    \item \textbf{CV‑трекинг} тест‑драйвов: YOLOv5, мультипоточность, кропы; точность 65\%→87\%, до 15 fps на Consumer GPU.
  \end{highlights}
\end{twocolentry}

% \vspace{0.5cm}
\section{Проекты}
\begin{twocolentry}{Сент. 2022 -- Май 2023}
  \textbf{НИУ ВШЭ НН} \textit{(Академический R\&D)}\newline
  ML Engineer | \textit{WhisperX, Llama.cpp, aiogram, guidance. \href{https://github.com/J1nsei/Voice-Message-Summarization-Bot}{GitHub}}
  \begin{highlights}
    \item Разработка автономного AI-ассистента (Telegram Bot на aiogram) для транскрибации и суммаризации аудио с фокусом на Data Privacy. Спроектировал ASR-пайплайн на базе WhisperX с выравниванием текста и regex-фильтрацией шума. Реализовал гибридный (CPU/GPU) инференс квантованной LLM через llama.cpp с кастомным управлением памятью и жестким контролем генерации через фреймворк guidance. Проект полностью независим от внешних API и оптимизирован для локального запуска на потребительском железе (от 12GB VRAM).
  \end{highlights}
\end{twocolentry}

%\vspace{0.5cm}
\begin{twocolentry}{Февр. 2021 -- Май 2021}
  \textbf{НИУ ВШЭ НН} \textit{(совместный R\&D проект с бизнесом)}\newline
  ML Researcher | \textit{Selenium, Colabeler, torchvision}
  \begin{highlights}
    \item Проектирование пайплайна сбора и подготовки данных для CV-модели детекции повреждений автомобилей. Сбор датасета, разметка, аугментация, баланс классов.
  \end{highlights}
\end{twocolentry}

\section{Образование}

\eduentry
{Национальный исследовательский университет «Высшая школа экономики»}
{Нижний Новгород, Россия}
{Магистратура, Интеллектуальный анализ данных}
{Сент. 2022 -- Июнь 2024}

\eduentry
{Национальный исследовательский университет «Высшая школа экономики»}
{Нижний Новгород, Россия}
{Бакалавриат, Прикладная математика и информатика}
{Сент. 2018 -- Июнь 2022}

\subsection{Дополнительно}
\subeduentry
{\href{https://education.yandex.ru/profile/docs?id=927802ab-4cde-4068-9354-99469527e3fc}{AI Agents Security Week}}
{27.07.2026 -- 31.07.2026}
\subeduentry
{\href{https://education.yandex.ru/profile/docs?id=2f1dd619-4a1b-45a6-93f1-491e18333322}{Agents Week Школы анализа данных (YSDA)}}
{06.04.2026 -- 10.04.2026}
\subeduentry
{\href{https://education.yandex.ru/profile/docs?id=95aa3c1d-031f-4a09-9d8e-af8159728780}{Интенсив ШАД Яндекс: LLM Scaling Week}}
{10.11.2025 -- 14.11.2025}

\section{Навыки}

\begin{onecolentry}
  \textbf{DL/ML:} Python, PyTorch, vLLM, transformers, LangChain, LangGraph, llama.cpp, deepeval, RAGFlow, NumPy, prompt engineering, infinity, unsloth, scikit-learn, WhisperX, spaCy.
\end{onecolentry}

\begin{onecolentry}
  \textbf{Infra/Tools:} Linux, Docker, FastAPI, Git, DVC, streamlit, LangFuse, W\&B, matplotlib, Grafana, Kubeflow
\end{onecolentry}

\begin{onecolentry}
  \textbf{Data:} pandas, polars, SQL
\end{onecolentry}

\begin{onecolentry}
  \textbf{Languages:} Russian - Native, English - B2 / Upper-Intermediate
\end{onecolentry}

\end{document}
